% Section 13. Drafted from docs/LANDSCAPE.md, with the corrections of docs/SUBSUMPTION.md.

\section{Related work}\label{sec:related}

\paragraph{CRDTs and coordination avoidance.}
CRDTs~\cite{shapiro2011} obtain convergence by designing operations that commute; \cref{sec:crdt}
shows they are the case where repair does nothing, that their conditions are exact on reachable states,
and that every convergent repaired system is a CRDT on its governed steps. The CALM
theorem~\cite{hellerstein2010,ameloot2013} characterizes the programs with a coordination-free
implementation as the monotone ones; the monotone-cycle regime of the map is related, a monotone
operator whose least fixed point is reached in any fair order, carrying business invariants.
Invariant confluence~\cite{bailis2014,whittaker2018} asks which transactions preserve an invariant
without coordination, with the invariant checked and never repaired (\cref{sec:crdt-iconf}).

\paragraph{Invariants on weakly consistent replication.}
RedBlue consistency~\cite{li2012} labels operations as commuting or globally ordered; Indigo and
IPA~\cite{balegas2015,balegas2018} combine invariant preservation with repair and with
reservations, building on escrow~\cite{oneil1986}. These systems ground convergence in operation
commutativity plus invariant preservation, transforming operations where needed. Normalization
confluence admits operations that do not commute and that violate the invariant, and proves
convergence by Newman's lemma on the repaired rewrite system: compensation commutativity is local
confluence, not operation commutativity. Bayou~\cite{terry1995} resolved conflicts by
application-specific merge procedures after a tentative order; sagas~\cite{garcia1987} use
compensation to undo partial work. Normalization confluence gives compensation a convergence theory:
exactly when repairs make concurrent orders agree.

\paragraph{Rewriting, fixed points and traces.}
The single-registry theorem is the critical-pair theorem~\cite{knuthbendix1970,huet1980} specialized to
a system whose critical pairs are normal forms, with Newman's lemma~\cite{newman1942} as the bridge.
Chaotic iteration to a least fixed point is classical in abstract interpretation~\cite{cousot1977} and
in discrete dynamics~\cite{robert1986}. Layer~A is trace theory~\cite{cartier1969,mazurkiewicz1977}
generalized to admissible words and non-commutation generators, with the congruence facts of semigroup
presentations~\cite{howie1995}.

\paragraph{Constraints and long-run behavior.}
Layer~B places consistency in the algebraic theory of constraint
satisfaction~\cite{jeavons1997,bulatov2017,zhuk2020}, and the identification of global sections with
constraint solutions follows the sheaf-theoretic treatment of contextuality~\cite{abramsky2011,
abramsky2014}. Layer~C's kernel is the fair-cycle characterization of model
checking~\cite{clarke1986,emersonlei1987,vardi1986}; its relation to fair
termination~\cite{lehmann1981,francez1986} is stated as open.

\paragraph{Mechanized convergence.}
Gomes et al.~\cite{gomes2017} verified strong eventual consistency for op-based CRDTs in Isabelle;
their theorem is recovered here in their own definitions (\cref{sec:tests-outside}). Dijkstra's
self-stabilizing ring has been verified in PVS~\cite{qadeer1998} and Coq~\cite{altisen2017}. The
development here differs in scope rather than in any single proof: it mechanizes an exact condition, a
hardness result or a listed open problem for every regime of a map that it does not exclude by design, with the counterexamples that fix
each hypothesis.
